% !TeX root = ../main.tex

\chapter{實驗結果與分析}

\section{實驗設計}

說明實驗設計...

\begin{table}[htbp]
  \centering
  \caption{實驗設定範例}
  \label{tab:experiment-setting}
  \begin{tabular}{lll}
    \toprule
    項目 & 設定 & 說明 \\
    \midrule
    資料集 & Dataset A & 可替換為實際資料來源 \\
    評估指標 & Accuracy & 可依研究領域調整 \\
    重複次數 & 5 & 用於估計平均表現 \\
    \bottomrule
  \end{tabular}
\end{table}

\section{實驗結果}

\subsection{結果一}

分析結果...

\begin{figure}[htbp]
  \centering
  \begin{tikzpicture}
    \draw[rounded corners, thick] (0,0) rectangle (3,1);
    \draw[rounded corners, thick] (4,0) rectangle (7,1);
    \draw[rounded corners, thick] (8,0) rectangle (11,1);
    \node at (1.5,0.5) {資料};
    \node at (5.5,0.5) {方法};
    \node at (9.5,0.5) {結果};
    \draw[->, thick] (3,0.5) -- (4,0.5);
    \draw[->, thick] (7,0.5) -- (8,0.5);
  \end{tikzpicture}
  \caption{研究流程示意圖}
  \label{fig:workflow}
\end{figure}

\section{討論}

深入討論結果含義...
